%% kapa: exam=1 round=26 year=2015 period=1 form=A subject=economics subject_ko=경제원론 seq=1 number=41
%% source: Q-Net 감정평가사 기출문제 (cst00309) artlSeq=5138001 file=extracted/1008/5138001/제26회_감정평가사 1차_기출/제26회_감정평가사 1차_기출/1교시_감평_A형.pdf
%% conversion: kordoc 4.24.51 -> tools/gen1.py (auto); review=proofread-sonnet; flags=tikz
\begin{question}{41}
다음과 같이 시장수요곡선(D)과 시장공급곡선(S)이 주어졌을 때, 정부가 생산자에게 세금을 부과하여 공급곡선이 S에서 S'로 이동하였다. 다음 중 옳은 것은? (단, 시장수요곡선은 완전 탄력적이며, 시장공급곡선은 우상향한다.)
\begin{center}
\begin{tikzpicture}[x=0.12mm,y=-0.12mm,font=\small,line width=0.9pt]
\draw[-{Stealth[length=1.6mm]}] (58,296) -- (58,15);
\draw[-{Stealth[length=1.6mm]}] (58,296) -- (333,296);
\draw (58,161) -- (323,161);
\draw (58,272) -- (292,98);
\draw[dashed] (58,212) -- (292,38);
\node at (343,160) {D};
\node at (316,90) {S};
\node at (320,22) {S$'$};
\node[align=center] at (28,48) {가\\격};
\node at (355,325) {수량};
\node at (35,312) {0};
\end{tikzpicture}
\end{center}
\begin{choices}
  \item 모든 세금은 소비자가 부담한다.
  \item 균형거래량은 변화가 없다.
  \item 생산자잉여는 감소한다.
  \item 소비자잉여는 증가한다.
  \item 정부의 조세수입은 발생하지 않는다.
\end{choices}
\end{question}
\answer{3}
